\newpage
\pagestyle{plain}
\chapter{Дүгнэлт ба санал зөвлөмж}

%% ============================================================
\section{Дүгнэлт}
%% ============================================================

Энэхүү төгсөлтийн судалгааны ажлын хүрээнд ``13-р зууны Монголын түүхийг таниулах систем'' нэртэй мобайл аппликейшнийг амжилттай хөгжүүлсэн. Аппыг Flutter кросс-платформ фреймворк ашиглан Dart програмчлалын хэл дээр бичсэн бөгөөд Android болон iOS үйлдлийн системд ажиллах боломжтой.

Төслийн гол үр дүнг дараах байдлаар дүгнэж болно:

\begin{enumerate}
	\item \textbf{7 үндсэн функцийг хэрэгжүүлсэн:} Түүхэн хүмүүсийн намтар, интерактив газрын зураг (3D/2D), цагийн хэлхээс, квиз, мультимедиа, гэр бүлийн мод, соёлын мэдээллийн дэлгэцүүдийг бүрэн хөгжүүлсэн.
	
	\item \textbf{Оффлайн эхлэлтэй архитектур:} JSON файлд суурилсан оффлайн өгөгдлийн систем бүтээж, сүлжээний холболтгүй орчинд ч бүрэн ажиллах боломжтой болгосон.
	
	\item \textbf{Provider төлөвийн удирдлага:} Flutter-ийн албан ёсны санал болгосон Provider загвар ашиглан цэвэр, модульчлагдсан архитектурыг хангасан.
	
	\item \textbf{3D интерактив газрын зураг:} \texttt{flutter\_earth\_globe} сан ашиглан Монголын эзэнт гүрний нутаг дэвсгэр, байлдан дагуулалтын маркерууд, аян замын чиглэлүүдийг 3D бөмбөрцөг дээр интерактив байдлаар харуулсан — энэ нь ижил төрлийн аппуудаас ялгарах гол онцлог юм.
	
	\item \textbf{Хэрэглэгчийн сэтгэл ханамж:} 15 хэрэглэгчийн ашиглалтын туршилтаар дундаж 4.4/5.0 оноо авсан нь аппын дизайн, контент, интерактив байдал хэрэглэгчдэд таалагдсаныг харуулж байна.
	
	\item \textbf{Firebase бэлтгэл:} Cloud синхрончлолын үндсэн дэд бүтцийг бэлтгэж, ирээдүйд онлайн функцуудыг нэмэхэд бэлэн болгосон.
\end{enumerate}

Судалгааны ажлын хүрээнд Agile аргачлалыг хэрэглэж, 8 сарын хугацаанд үе шат бүрийг (судалгаа, дизайн, хөгжүүлэлт, туршилт) амжилттай гүйцэтгэсэн. Модульчлагдсан кодын бүтэц нь цаашдын өргөтгөл, засварлалтыг хялбар болгосон.

%% ============================================================
\section{Санал зөвлөмж}
%% ============================================================

Цаашид аппыг дараах чиглэлүүдээр сайжруулах, өргөтгөх боломжтой:

\subsection{Богино хугацааны сайжруулалт}

\begin{enumerate}
	\item \textbf{Firebase Cloud Firestore интеграци:} Бэлтгэсэн Firebase дэд бүтцийг бүрэн идэвхжүүлж, шинэ контентыг Cloud-оос татах, локал кэштэй синхрончлох функцийг нэмэх.
	\item \textbf{SQLite өгөгдлийн сан:} JSON файлын оронд SQLite ашиглан өгөгдлийн удирдлагыг сайжруулах. Энэ нь хайлт, шүүлтүүрийн хурдыг нэмэгдүүлнэ.
	\item \textbf{Квизийн системийг сайжруулах:} Түвшин (хялбар, дунд, хэцүү) нэмэх, Rich Media (зурагтай) асуулт, хугацааны хязгаарлалт зэргийг хэрэгжүүлэх.
	\item \textbf{Аудио тайлбар:} Намтар, үйл явдлуудад аудио тайлбар нэмж, хэрэглэгчийн туршлагыг баяжуулах.
	\item \textbf{iOS туршилт:} Apple-ийн төхөөрөмж дээр бүрэн туршилт хийж, App Store-д байршуулах.
\end{enumerate}

\subsection{Урт хугацааны өргөтгөл}

\begin{enumerate}
	\item \textbf{Хэл дэмжлэг:} Англи, Хятад, Орос зэрэг хэл дээрх орчуулга нэмж, олон улсын хэрэглэгчдэд хүрэх.
	\item \textbf{AR/VR элемент:} Нэмэгдсэн бодит байдлын (Augmented Reality) технологи ашиглан түүхэн газруудыг 3D-ээр сэргээн харуулах.
	\item \textbf{Хиймэл оюун ухаан:} ML Kit эсвэл GPT API ашиглан хэрэглэгчид зориулсан хувийн санал болголт (recommendation) системийг хөгжүүлэх.
	\item \textbf{Нийгмийн функц:} Хэрэглэгчид хоорондоо мэдлэг хуваалцах, квизэд өрсөлдөх боломжийг нэмэх.
	\item \textbf{Түүхийн цар хүрээг өргөтгөх:} 13-р зуунаас гадна Монголын бүх түүхийг хамарсан иж бүрэн боловсролын платформ болгох.
	\item \textbf{Багш нарт зориулсан модуль:} Хичээлийн хөтөлбөртэй уялдуулсан багшийн удирдлагын самбар (dashboard) бүтээх.
\end{enumerate}

\subsection{Ерөнхий зөвлөмж}

Монгол Улсад боловсролын зорилготой мобайл аппликейшний хэрэгцээ өндөр байгаа бөгөөд ялангуяа Монгол хэл дээрх чанартай, интерактив контент дутагдалтай байна. Энэхүү судалгааны ажлын үр дүн нь дижитал боловсролын салбарт хувь нэмэр оруулж, Монголын түүхийг залуу үед хүртээмжтэй хэлбэрээр хүргэх боломжийг нэмэгдүүлнэ гэж үзэж байна.

Цаашдын судлаачдад дараах зүйлийг санал болгож байна:
\begin{itemize}
	\item Хэрэглэгчийн судалгааг өргөн хүрээнд хийж, контентыг байнга шинэчлэх
	\item Боловсролын байгууллагуудтай хамтран ажиллаж, хичээлийн хөтөлбөртэй уялдуулах
	\item Аппын гүйцэтгэлийг (performance) байнга хянаж, оновчлох
	\item Open source болгож, олон хөгжүүлэгчдийн оролцоог дэмжих
\end{itemize}
